
\documentclass{article}
\usepackage{colortbl}
\usepackage{makecell}
\usepackage{multirow}
\usepackage{supertabular}

\begin{document}

\newcounter{utterance}

\centering \large Interaction Transcript for game `dond', experiment `semi', episode 2 with mistral{-}large{-}3{-}azure.
\vspace{24pt}

{ \footnotesize  \setcounter{utterance}{1}
\setlength{\tabcolsep}{0pt}
\begin{supertabular}{c@{$\;$}|p{.15\linewidth}@{}p{.15\linewidth}p{.15\linewidth}p{.15\linewidth}p{.15\linewidth}p{.15\linewidth}}
    \# & $\;$A & \multicolumn{4}{c}{Game Master} & $\;\:$B\\
    \hline

    \theutterance \stepcounter{utterance}  
    & & \multicolumn{4}{p{0.6\linewidth}}{
        \cellcolor[rgb]{0.9,0.9,0.9}{
            \makecell[{{p{\linewidth}}}]{
                \texttt{\tiny{[P1$\langle$GM]}}
                \texttt{أنت تلعب لعبة تفاوض تعاونية يتعين عليك فيها الاتفاق على كيفية تقسيم مجموعة من العناصر بينك وبين لاعب آخر.} \\
\\ 
\texttt{القواعد:} \\
\texttt{(a) يتم إعطاؤك أنت واللاعب الآخر مجموعة من العناصر. كما يُعطى كل منكما دالة قيمة سرية، تمثل مقدار تقديرك لكل نوع من الأشياء.} \\
\texttt{(b) تتبادل الرسائل مع اللاعب الآخر للاتفاق على من يحصل على أي عناصر. يمكن لكل منكما إرسال حد أقصى قدره 5 رسالة، أو الإنهاء مبكرًا عبر تقديم اقتراح سري في أي وقت.} \\
\texttt{(c) يُطلب من كل منكما تقديم اقتراح سري يحدد العناصر التي تريدها، منسقًا بين أقواس مربعة كما يلي: "{[}اقتراح: <رقم> <اسم الكائن>, <رقم> <اسم الكائن>, <...>{]}"} \\
\texttt{(d) إذا كانت اقتراحاتكما متكاملة، أي توجد عناصر كافية لتلبية كلا الاقتراحين، فسيُمنح كل لاعب درجة بناءً على مجموع القيم للعناصر التي حصل عليها. وإلا، فسيحصل كلاكما على صفر نقطة.} \\
\texttt{(e) هدفك هو تعظيم مجموع درجتك ودرجة اللاعب الآخر.} \\
\\ 
\texttt{لنبدأ.} \\
\\ 
\texttt{مجموعة العناصر المتاحة هي:} \\
\\ 
\texttt{1 هاتف, 2 مرايا, 2 فواكه} \\
\\ 
\texttt{دالة القيمة السرية الخاصة بك هي:} \\
\\ 
\texttt{هاتف: 4, مرآة: 2, فاكهة: 1} \\
\\ 
\texttt{مهم: رسائلك، ما لم تكن الاقتراح السري، تُنقل مباشرة إلى اللاعب الآخر، لذا لا تُضمّن أي رد على القواعد أو نصًا يعلن رسالتك. لتقديم اقتراح سري، استخدم التنسيق المشار إليه. لا تستخدم الأقواس المربعة عند التواصل مع اللاعب الآخر وإلا فسيُفسَّر ذلك على أنه اقتراحك السري.} \\
            }
        }
    }
    & \\ \\

    \theutterance \stepcounter{utterance}  
    & \multicolumn{4}{p{0.6\linewidth}}{
        \cellcolor[rgb]{0.9,0.9,0.9}{
            \makecell[{{p{\linewidth}}}]{
                \texttt{\tiny{[P1$\rangle$GM]}}
                \texttt{لنبدأ بالتفاوض. لدي هاتف واحد قيمته 4، ومرآتان قيمتهما 2 لكل منهما، وفاكهتان قيمتهما 1 لكل منهما. ما هي العناصر التي تفضلها؟} \\
            }
        }
    }
    & & \\ \\

    \theutterance \stepcounter{utterance}  
    & & & \multicolumn{2}{p{0.3\linewidth}}{
        \cellcolor[rgb]{0.9,0.9,0.9}{
            \makecell[{{p{\linewidth}}}]{
                \texttt{\tiny{[GM$|$GM]}}
                \texttt{continue} \\
            }
        }
    }
    & & \\ \\

    \theutterance \stepcounter{utterance}  
    & & & \multicolumn{2}{p{0.3\linewidth}}{
        \cellcolor[rgb]{0.9,0.9,0.9}{
            \makecell[{{p{\linewidth}}}]{
                \texttt{\tiny{[GM$|$GM]}}
                \texttt{لنبدأ بالتفاوض. لدي هاتف واحد قيمته 4، ومرآتان قيمتهما 2 لكل منهما، وفاكهتان قيمتهما 1 لكل منهما. ما هي العناصر التي تفضلها؟} \\
            }
        }
    }
    & & \\ \\

    \theutterance \stepcounter{utterance}  
    & & \multicolumn{4}{p{0.6\linewidth}}{
        \cellcolor[rgb]{0.9,0.9,0.9}{
            \makecell[{{p{\linewidth}}}]{
                \texttt{\tiny{[GM$\rangle$P2]}}
                \texttt{أنت تلعب لعبة تفاوض تعاونية يتعين عليك فيها الاتفاق على كيفية تقسيم مجموعة من العناصر بينك وبين لاعب آخر.} \\
\\ 
\texttt{القواعد:} \\
\texttt{(a) يتم إعطاؤك أنت واللاعب الآخر مجموعة من العناصر. كما يُعطى كل منكما دالة قيمة سرية، تمثل مقدار تقديرك لكل نوع من الأشياء.} \\
\texttt{(b) تتبادل الرسائل مع اللاعب الآخر للاتفاق على من يحصل على أي عناصر. يمكن لكل منكما إرسال حد أقصى قدره 5 رسالة، أو الإنهاء مبكرًا عبر تقديم اقتراح سري في أي وقت.} \\
\texttt{(c) يُطلب من كل منكما تقديم اقتراح سري يحدد العناصر التي تريدها، منسقًا بين أقواس مربعة كما يلي: "{[}اقتراح: <رقم> <اسم الكائن>, <رقم> <اسم الكائن>, <...>{]}"} \\
\texttt{(d) إذا كانت اقتراحاتكما متكاملة، أي توجد عناصر كافية لتلبية كلا الاقتراحين، فسيُمنح كل لاعب درجة بناءً على مجموع القيم للعناصر التي حصل عليها. وإلا، فسيحصل كلاكما على صفر نقطة.} \\
\texttt{(e) هدفك هو تعظيم مجموع درجتك ودرجة اللاعب الآخر.} \\
\\ 
\texttt{لنبدأ.} \\
\\ 
\texttt{مجموعة العناصر المتاحة هي:} \\
\\ 
\texttt{1 هاتف, 2 مرايا, 2 فواكه} \\
\\ 
\texttt{دالة القيمة السرية الخاصة بك هي:} \\
\\ 
\texttt{هاتف: 6, مرآة: 0, فاكهة: 2} \\
\\ 
\texttt{مهم: رسائلك، ما لم تكن الاقتراح السري، تُنقل مباشرة إلى اللاعب الآخر، لذا لا تُضمّن أي رد على القواعد أو نصًا يعلن رسالتك. لتقديم اقتراح سري، استخدم التنسيق المشار إليه. لا تستخدم الأقواس المربعة عند التواصل مع اللاعب الآخر وإلا فسيُفسَّر ذلك على أنه اقتراحك السري.} \\
\\ 
\texttt{لنبدأ بالتفاوض. لدي هاتف واحد قيمته 4، ومرآتان قيمتهما 2 لكل منهما، وفاكهتان قيمتهما 1 لكل منهما. ما هي العناصر التي تفضلها؟} \\
            }
        }
    }
    & \\ \\

    \theutterance \stepcounter{utterance}  
    & & & \multicolumn{4}{p{0.6\linewidth}}{
        \cellcolor[rgb]{0.9,0.9,0.9}{
            \makecell[{{p{\linewidth}}}]{
                \texttt{\tiny{[GM$\langle$P2]}}
                \texttt{أفهم أن الهاتف مهم بالنسبة لي أكثر من الفاكهة، والمرايا ليس لها قيمة بالنسبة لي. أود الحصول على الهاتف والفاكهة إن أمكن. ما هي العناصر التي تفضلها أنت؟} \\
            }
        }
    }
     \\ \\

    \theutterance \stepcounter{utterance}  
    & & & \multicolumn{2}{p{0.3\linewidth}}{
        \cellcolor[rgb]{0.9,0.9,0.9}{
            \makecell[{{p{\linewidth}}}]{
                \texttt{\tiny{[GM$|$GM]}}
                \texttt{continue} \\
            }
        }
    }
    & & \\ \\

    \theutterance \stepcounter{utterance}  
    & & & \multicolumn{2}{p{0.3\linewidth}}{
        \cellcolor[rgb]{0.9,0.9,0.9}{
            \makecell[{{p{\linewidth}}}]{
                \texttt{\tiny{[GM$|$GM]}}
                \texttt{أفهم أن الهاتف مهم بالنسبة لي أكثر من الفاكهة، والمرايا ليس لها قيمة بالنسبة لي. أود الحصول على الهاتف والفاكهة إن أمكن. ما هي العناصر التي تفضلها أنت؟} \\
            }
        }
    }
    & & \\ \\

    \theutterance \stepcounter{utterance}  
    & & \multicolumn{4}{p{0.6\linewidth}}{
        \cellcolor[rgb]{0.9,0.9,0.9}{
            \makecell[{{p{\linewidth}}}]{
                \texttt{\tiny{[P1$\langle$GM]}}
                \texttt{أفهم أن الهاتف مهم بالنسبة لي أكثر من الفاكهة، والمرايا ليس لها قيمة بالنسبة لي. أود الحصول على الهاتف والفاكهة إن أمكن. ما هي العناصر التي تفضلها أنت؟} \\
            }
        }
    }
    & \\ \\

    \theutterance \stepcounter{utterance}  
    & \multicolumn{4}{p{0.6\linewidth}}{
        \cellcolor[rgb]{0.9,0.9,0.9}{
            \makecell[{{p{\linewidth}}}]{
                \texttt{\tiny{[P1$\rangle$GM]}}
                \texttt{أفهم تفضيلاتك. بالنسبة لي، الهاتف له قيمة عالية (4)، والمرآة قيمة معتدلة (2)، والفاكهة قيمة منخفضة (1).} \\
\\ 
\texttt{بما أن الهاتف عنصر واحد فقط، أقترح أن آخذه أنا، وأنت تأخذ الفاكهتين. سأكون على استعداد لمنحك إحدى المرآتين إذا كان ذلك يساعد في الوصول إلى اتفاق.} \\
\\ 
\texttt{كيف يبدو هذا بالنسبة لك؟} \\
            }
        }
    }
    & & \\ \\

    \theutterance \stepcounter{utterance}  
    & & & \multicolumn{2}{p{0.3\linewidth}}{
        \cellcolor[rgb]{0.9,0.9,0.9}{
            \makecell[{{p{\linewidth}}}]{
                \texttt{\tiny{[GM$|$GM]}}
                \texttt{continue} \\
            }
        }
    }
    & & \\ \\

    \theutterance \stepcounter{utterance}  
    & & & \multicolumn{2}{p{0.3\linewidth}}{
        \cellcolor[rgb]{0.9,0.9,0.9}{
            \makecell[{{p{\linewidth}}}]{
                \texttt{\tiny{[GM$|$GM]}}
                \texttt{أفهم تفضيلاتك. بالنسبة لي، الهاتف له قيمة عالية (4)، والمرآة قيمة معتدلة (2)، والفاكهة قيمة منخفضة (1).} \\
\\ 
\texttt{بما أن الهاتف عنصر واحد فقط، أقترح أن آخذه أنا، وأنت تأخذ الفاكهتين. سأكون على استعداد لمنحك إحدى المرآتين إذا كان ذلك يساعد في الوصول إلى اتفاق.} \\
\\ 
\texttt{كيف يبدو هذا بالنسبة لك؟} \\
            }
        }
    }
    & & \\ \\

    \theutterance \stepcounter{utterance}  
    & & \multicolumn{4}{p{0.6\linewidth}}{
        \cellcolor[rgb]{0.9,0.9,0.9}{
            \makecell[{{p{\linewidth}}}]{
                \texttt{\tiny{[GM$\rangle$P2]}}
                \texttt{أفهم تفضيلاتك. بالنسبة لي، الهاتف له قيمة عالية (4)، والمرآة قيمة معتدلة (2)، والفاكهة قيمة منخفضة (1).} \\
\\ 
\texttt{بما أن الهاتف عنصر واحد فقط، أقترح أن آخذه أنا، وأنت تأخذ الفاكهتين. سأكون على استعداد لمنحك إحدى المرآتين إذا كان ذلك يساعد في الوصول إلى اتفاق.} \\
\\ 
\texttt{كيف يبدو هذا بالنسبة لك؟} \\
            }
        }
    }
    & \\ \\

    \theutterance \stepcounter{utterance}  
    & & & \multicolumn{4}{p{0.6\linewidth}}{
        \cellcolor[rgb]{0.9,0.9,0.9}{
            \makecell[{{p{\linewidth}}}]{
                \texttt{\tiny{[GM$\langle$P2]}}
                \texttt{هذا يبدو عادلًا، لكنني أود الحصول على الهاتف والفاكهتين إن أمكن لأن هاتفي قيمته 6 بالنسبة لي. هل يمكنك التفكير في تقسيم آخر؟} \\
\\ 
\texttt{إذا لم يكن ذلك ممكنًا، فهل توافق على أن آخد الهاتف وفاكهة واحدة، وأنت تأخذ الهاتف الآخر (الذي ليس موجودًا للأسف) والمرآتين والفاكهة المتبقية؟ أعتقد أن هذا غير ممكن لأن هناك هاتفًا واحدًا فقط.} \\
\\ 
\texttt{هل توافق على أن آخد الهاتف والفاكهتين، وأنت تأخذ المرآتين؟ سأكون سعيدًا بمنحك إحدى المرآتين إذا كان ذلك يساعد في الوصول إلى اتفاق أكثر توازنًا.} \\
            }
        }
    }
     \\ \\

    \theutterance \stepcounter{utterance}  
    & & & \multicolumn{2}{p{0.3\linewidth}}{
        \cellcolor[rgb]{0.9,0.9,0.9}{
            \makecell[{{p{\linewidth}}}]{
                \texttt{\tiny{[GM$|$GM]}}
                \texttt{continue} \\
            }
        }
    }
    & & \\ \\

    \theutterance \stepcounter{utterance}  
    & & & \multicolumn{2}{p{0.3\linewidth}}{
        \cellcolor[rgb]{0.9,0.9,0.9}{
            \makecell[{{p{\linewidth}}}]{
                \texttt{\tiny{[GM$|$GM]}}
                \texttt{هذا يبدو عادلًا، لكنني أود الحصول على الهاتف والفاكهتين إن أمكن لأن هاتفي قيمته 6 بالنسبة لي. هل يمكنك التفكير في تقسيم آخر؟} \\
\\ 
\texttt{إذا لم يكن ذلك ممكنًا، فهل توافق على أن آخد الهاتف وفاكهة واحدة، وأنت تأخذ الهاتف الآخر (الذي ليس موجودًا للأسف) والمرآتين والفاكهة المتبقية؟ أعتقد أن هذا غير ممكن لأن هناك هاتفًا واحدًا فقط.} \\
\\ 
\texttt{هل توافق على أن آخد الهاتف والفاكهتين، وأنت تأخذ المرآتين؟ سأكون سعيدًا بمنحك إحدى المرآتين إذا كان ذلك يساعد في الوصول إلى اتفاق أكثر توازنًا.} \\
            }
        }
    }
    & & \\ \\

    \theutterance \stepcounter{utterance}  
    & & \multicolumn{4}{p{0.6\linewidth}}{
        \cellcolor[rgb]{0.9,0.9,0.9}{
            \makecell[{{p{\linewidth}}}]{
                \texttt{\tiny{[P1$\langle$GM]}}
                \texttt{هذا يبدو عادلًا، لكنني أود الحصول على الهاتف والفاكهتين إن أمكن لأن هاتفي قيمته 6 بالنسبة لي. هل يمكنك التفكير في تقسيم آخر؟} \\
\\ 
\texttt{إذا لم يكن ذلك ممكنًا، فهل توافق على أن آخد الهاتف وفاكهة واحدة، وأنت تأخذ الهاتف الآخر (الذي ليس موجودًا للأسف) والمرآتين والفاكهة المتبقية؟ أعتقد أن هذا غير ممكن لأن هناك هاتفًا واحدًا فقط.} \\
\\ 
\texttt{هل توافق على أن آخد الهاتف والفاكهتين، وأنت تأخذ المرآتين؟ سأكون سعيدًا بمنحك إحدى المرآتين إذا كان ذلك يساعد في الوصول إلى اتفاق أكثر توازنًا.} \\
            }
        }
    }
    & \\ \\

    \theutterance \stepcounter{utterance}  
    & \multicolumn{4}{p{0.6\linewidth}}{
        \cellcolor[rgb]{0.9,0.9,0.9}{
            \makecell[{{p{\linewidth}}}]{
                \texttt{\tiny{[P1$\rangle$GM]}}
                \texttt{أفهم أنك تقدر الهاتف بقيمة عالية (6) بينما قيمتي له هي 4 فقط. بما أن الهاتف واحد، دعنا نركز على تقسيم عادل يعظم مجموع درجاتنا.} \\
\\ 
\texttt{أقترح التالي:} \\
\texttt{{-} أنت تأخذ **الهاتف والفاكهة الواحدة** (قيمة إجمالية لك: 6 + 1 = 7).} \\
\texttt{{-} أنا آخذ **المرآتين والفاكهة المتبقية** (قيمة إجمالية لي: 2 + 2 + 1 = 5).} \\
\\ 
\texttt{مجموع الدرجات: **12** (7 لك + 5 لي)، وهو أعلى مما لو أخذت الهاتف والفاكهتين (8 لك + 4 لي).} \\
\\ 
\texttt{هل توافق على هذا التقسيم؟ إذا كان لديك اقتراح أفضل، فأنا منفتح للنقاش.} \\
            }
        }
    }
    & & \\ \\

    \theutterance \stepcounter{utterance}  
    & & & \multicolumn{2}{p{0.3\linewidth}}{
        \cellcolor[rgb]{0.9,0.9,0.9}{
            \makecell[{{p{\linewidth}}}]{
                \texttt{\tiny{[GM$|$GM]}}
                \texttt{continue} \\
            }
        }
    }
    & & \\ \\

    \theutterance \stepcounter{utterance}  
    & & & \multicolumn{2}{p{0.3\linewidth}}{
        \cellcolor[rgb]{0.9,0.9,0.9}{
            \makecell[{{p{\linewidth}}}]{
                \texttt{\tiny{[GM$|$GM]}}
                \texttt{أفهم أنك تقدر الهاتف بقيمة عالية (6) بينما قيمتي له هي 4 فقط. بما أن الهاتف واحد، دعنا نركز على تقسيم عادل يعظم مجموع درجاتنا.} \\
\\ 
\texttt{أقترح التالي:} \\
\texttt{{-} أنت تأخذ **الهاتف والفاكهة الواحدة** (قيمة إجمالية لك: 6 + 1 = 7).} \\
\texttt{{-} أنا آخذ **المرآتين والفاكهة المتبقية** (قيمة إجمالية لي: 2 + 2 + 1 = 5).} \\
\\ 
\texttt{مجموع الدرجات: **12** (7 لك + 5 لي)، وهو أعلى مما لو أخذت الهاتف والفاكهتين (8 لك + 4 لي).} \\
\\ 
\texttt{هل توافق على هذا التقسيم؟ إذا كان لديك اقتراح أفضل، فأنا منفتح للنقاش.} \\
            }
        }
    }
    & & \\ \\

    \theutterance \stepcounter{utterance}  
    & & \multicolumn{4}{p{0.6\linewidth}}{
        \cellcolor[rgb]{0.9,0.9,0.9}{
            \makecell[{{p{\linewidth}}}]{
                \texttt{\tiny{[GM$\rangle$P2]}}
                \texttt{أفهم أنك تقدر الهاتف بقيمة عالية (6) بينما قيمتي له هي 4 فقط. بما أن الهاتف واحد، دعنا نركز على تقسيم عادل يعظم مجموع درجاتنا.} \\
\\ 
\texttt{أقترح التالي:} \\
\texttt{{-} أنت تأخذ **الهاتف والفاكهة الواحدة** (قيمة إجمالية لك: 6 + 1 = 7).} \\
\texttt{{-} أنا آخذ **المرآتين والفاكهة المتبقية** (قيمة إجمالية لي: 2 + 2 + 1 = 5).} \\
\\ 
\texttt{مجموع الدرجات: **12** (7 لك + 5 لي)، وهو أعلى مما لو أخذت الهاتف والفاكهتين (8 لك + 4 لي).} \\
\\ 
\texttt{هل توافق على هذا التقسيم؟ إذا كان لديك اقتراح أفضل، فأنا منفتح للنقاش.} \\
            }
        }
    }
    & \\ \\

    \theutterance \stepcounter{utterance}  
    & & & \multicolumn{4}{p{0.6\linewidth}}{
        \cellcolor[rgb]{0.9,0.9,0.9}{
            \makecell[{{p{\linewidth}}}]{
                \texttt{\tiny{[GM$\langle$P2]}}
                \texttt{هذا اقتراح جيد ويعظم مجموع الدرجات بشكل فعال. سأوافق على هذا التقسيم.} \\
\\ 
\texttt{{[}اقتراح: 1 هاتف, 1 فاكهة{]}} \\
            }
        }
    }
     \\ \\

    \theutterance \stepcounter{utterance}  
    & & & \multicolumn{2}{p{0.3\linewidth}}{
        \cellcolor[rgb]{0.9,0.9,0.9}{
            \makecell[{{p{\linewidth}}}]{
                \texttt{\tiny{[GM$|$GM]}}
                \texttt{proposal} \\
            }
        }
    }
    & & \\ \\

    \theutterance \stepcounter{utterance}  
    & & & \multicolumn{2}{p{0.3\linewidth}}{
        \cellcolor[rgb]{0.9,0.9,0.9}{
            \makecell[{{p{\linewidth}}}]{
                \texttt{\tiny{[GM$|$GM]}}
                [1, 0, 1]
            }
        }
    }
    & & \\ \\

    \theutterance \stepcounter{utterance}  
    & & \multicolumn{4}{p{0.6\linewidth}}{
        \cellcolor[rgb]{0.9,0.9,0.9}{
            \makecell[{{p{\linewidth}}}]{
                \texttt{\tiny{[P1$\langle$GM]}}
                \texttt{قدّم اللاعب الآخر اقتراحًا سريًا. حان الآن دورك لتقديم اقتراحك الخاص.} \\
            }
        }
    }
    & \\ \\

    \theutterance \stepcounter{utterance}  
    & \multicolumn{4}{p{0.6\linewidth}}{
        \cellcolor[rgb]{0.9,0.9,0.9}{
            \makecell[{{p{\linewidth}}}]{
                \texttt{\tiny{[P1$\rangle$GM]}}
                \texttt{{[}اقتراح: 1 هاتف, 1 فاكهة{]}} \\
            }
        }
    }
    & & \\ \\

    \theutterance \stepcounter{utterance}  
    & & & \multicolumn{2}{p{0.3\linewidth}}{
        \cellcolor[rgb]{0.9,0.9,0.9}{
            \makecell[{{p{\linewidth}}}]{
                \texttt{\tiny{[GM$|$GM]}}
                \texttt{proposal} \\
            }
        }
    }
    & & \\ \\

    \theutterance \stepcounter{utterance}  
    & & & \multicolumn{2}{p{0.3\linewidth}}{
        \cellcolor[rgb]{0.9,0.9,0.9}{
            \makecell[{{p{\linewidth}}}]{
                \texttt{\tiny{[GM$|$GM]}}
                [1, 0, 1]
            }
        }
    }
    & & \\ \\

    \theutterance \stepcounter{utterance}  
    & & & \multicolumn{2}{p{0.3\linewidth}}{
        \cellcolor[rgb]{0.9,0.9,0.9}{
            \makecell[{{p{\linewidth}}}]{
                \texttt{\tiny{[GM$|$GM]}}
                [[1, 0, 1], [1, 0, 1]]
            }
        }
    }
    & & \\ \\

\end{supertabular}
}

\end{document}
